\documentclass[11pt]{article}
\usepackage{amsmath,amssymb,amsthm}
\usepackage[margin=1in]{geometry}
\usepackage{hyperref}
\usepackage{booktabs}

\newtheorem{theorem}{Theorem}
\newtheorem{lemma}[theorem]{Lemma}
\newtheorem{corollary}[theorem]{Corollary}
\theoremstyle{remark}
\newtheorem{remark}[theorem]{Remark}

\newcommand{\binomtwo}[1]{\binom{#1}{2}}

\title{Binomial coefficients $\binom{n}{2}$ that are products of consecutive primes}
\author{Srihari Mysore}
\date{DRAFT --- \today}

\begin{document}
\maketitle

\begin{abstract}
Erd\H{o}s and Graham asked whether $\binom{n}{k}$ can be a product of consecutive primes
infinitely often (Erd\H{o}s Problem \#386). For $k=2$ the known solutions are
$n = 4, 6, 15, 21, 715$. We prove an elementary \emph{balance inequality}: if
$\binom{n}{2} = p_a p_{a+1}\cdots p_{a+k-1}$ with $p_a \ge 5$, then
$(p_{a+k-1}/p_a)^{\lfloor k/2\rfloor} \ge 2 - 1/p_a$. Combined with explicit
short-interval results for primes, this reduces the problem to a finite computation for
every bound on $n$. We show there is no further solution with $n \le 10^{682}$, and with
$n \le 10^{9614}$ assuming a result of Dusart (2018). Unconditionally in $n$, any further
solution must be a product of at least $137$ consecutive primes (at least $1924$ under the
same result). An exhaustive search shows that no block starting at a prime below $100$ and
ending at a prime below $4\cdot 10^9$ gives a new solution. The core arithmetic inequality
is formalized in Lean~4.
\end{abstract}

\section{Introduction}

Let $p_0 = 2, p_1 = 3, \dots$ denote the primes. A \emph{block} is a product
$B(a,k) = p_a p_{a+1} \cdots p_{a+k-1}$ of $k \ge 1$ consecutive primes. Erd\H{o}s and
Graham \cite{EG} asked whether a binomial coefficient $\binom{n}{k}$, $2 \le k \le n-2$, can
equal a block infinitely often. For $k=2$ they remark that a proof that this happens only
finitely often ``seems hopeless''. The problem appears as Problem \#386 at
\cite{EP}, and the case $k = 2$ is formalized as \texttt{erdos\_386.variants.two} in
\cite{FC}.

For $k=2$ the question is: for which $n \ge 4$ is $n(n-1)/2$ a block? The known solutions
are
\[
\binomtwo{4} = 2\cdot 3,\quad \binomtwo{6} = 3\cdot 5,\quad \binomtwo{15} = 3\cdot5\cdot7,\quad
\binomtwo{21} = 2\cdot3\cdot5\cdot7,\quad \binomtwo{715} = 3\cdot5\cdot7\cdot11\cdot13\cdot17.
\]
The OEIS entry A280992 \cite{OEIS} reports no further solution with $n \le 5\cdot10^6$.
Blocks starting at $3$ correspond to $n(n-1) = r\#$ (a primorial). This is the question of
Nelson, Penney and Pomerance \cite{NPP}, who searched the first $3049$ primes.

\paragraph{Results.} Our main results are:
\begin{enumerate}
\item (Theorem~\ref{thm:bounded}) There is no solution other than the five above with
  $n \le 10^{12}$ (unconditionally), with $n \le 10^{682}$ (using Dusart 2010), and with
  $n \le 10^{9614}$ (using Dusart 2018).
\item (Theorem~\ref{thm:length}) For all $n$: any further solution is a block of
  $k \ge 137$ primes (Dusart 2010), and of $k \ge 1924$ primes (Dusart 2018).
\item (Theorem~\ref{thm:small}) No block $B(a,k)$ with $p_a < 100$ and
  $p_{a+k-1} \le 4\cdot10^9$ gives a new solution.
\end{enumerate}
All three rest on the balance inequality (Lemma~\ref{lem:balance}), whose proof is
elementary.

\section{The balance inequality}

\begin{lemma}\label{lem:notprime}
For $n \ge 4$, $\binomtwo{n}$ is not prime.
\end{lemma}
\begin{proof}
$\binomtwo{n}$ equals $(n/2)(n-1)$ or $n\cdot\frac{n-1}{2}$, and both factors are at least $2$.
\end{proof}

\begin{lemma}[Balance inequality]\label{lem:balance}
Let $n \ge 4$ and suppose $\binomtwo{n} = B(a,k)$ with $p = p_a \ge 5$. Let $r = p_{a+k-1}$
and $h = \lfloor k/2 \rfloor$. Then
\[
p\, r^{h} \;\ge\; (2p-1)\, p^{h}, \qquad\text{equivalently}\qquad
\Big(\frac{r}{p}\Big)^{h} \ge 2 - \frac1p .
\]
\end{lemma}

\begin{proof}
Let $E$ be the even and $O$ the odd element of $\{n, n-1\}$, and put $A = E/2$. Then
$\binomtwo{n} = AO$ with $\gcd(A,O) = 1$ and $|2A - O| = 1$. Every prime in the block is
$\ge 5$, so $AO$ is odd and squarefree. The block's primes therefore split as $S \sqcup T$
with $A = \prod S$ and $O = \prod T$. Since $n \ge 4$ we have $A \ge 2$ and $O \ge 3$, so
$S$ and $T$ are non-empty. Put $s = |S|$ and $t = |T|$, so $s+t=k$, and note that every
element of $S \cup T$ lies in $[p, r]$.

\emph{Case $s=t$.} Then $O/A = 2 \pm 1/A \ge 2 - 1/A \ge 2 - 1/p$, while
$O/A \le r^s/p^s$.

\emph{Case $s<t$.} Then $p^{s+1} \le p^t \le O \le 2A+1 \le 2r^s+1$, so
$(r/p)^s \ge p/2 - 1/(2p^s) \ge p/2 - 1/(2p) \ge 2 - 1/p$, since $p^2 - 4p + 1 \ge 0$.

\emph{Case $s>t$.} Then $2p^{t+1} \le 2A \le O+1 \le r^t+1$, so
$(r/p)^t \ge 2p - p^{-t} \ge 2 - 1/p$.

In every case $(r/p)^{\min(s,t)} \ge 2 - 1/p$. Since $\min(s,t) \le h$ and $r \ge p$, the
claim follows.
\end{proof}

\begin{remark}
The inequality is sharp: $\binomtwo{14} = 91 = 7\cdot 13$ has $p=7$, $r=13$, $h=1$, and
$13/7 = 2 - 1/7$. (This is not a solution, because $7$ and $13$ are not consecutive
primes.) The proof does not use consecutiveness. We checked the general form numerically
for all $n \le 3\cdot 10^6$.
\end{remark}

For a fixed start $a$, the left-hand side of Lemma~\ref{lem:balance} is non-decreasing in
$k$. Hence:

\begin{corollary}[Monotonicity]\label{cor:mono}
If the inequality of Lemma~\ref{lem:balance} fails for the block $B(a,k_0)$, it fails for
every $B(a,k)$ with $k \le k_0$.
\end{corollary}

\begin{lemma}[Integer test]\label{lem:test}
Let $p \ge 41$, $d = r-p \ge 0$ and $h \ge 0$. If $100\,h\,d < 68\,p$, then
$(r/p)^h < 2 - 1/p$.
\end{lemma}
\begin{proof}
$(1+d/p)^h \le e^{hd/p} < e^{0.68} < 1.974 < 2 - 1/41 \le 2 - 1/p$.
\end{proof}

\section{Large starting primes}

We use explicit short-interval theorems of the following form:
\begin{quote}
(G) For all $x \ge x_0$ there is a prime in $(x, x(1+\varepsilon(x))]$, where
$\varepsilon$ is non-increasing.
\end{quote}
Dusart \cite[Prop.~6.8]{D10} proves (G) with $x_0 = 396738$ and
$\varepsilon(x) = 1/(25\ln^2 x)$. Dusart \cite[Cor.~5.5]{D18} (see also \cite{Ax}) proves it
with $x_0 = 468991632$ and $\varepsilon(x) = 1/(5000\ln^2 x)$. Let
$X_0 = 10\,726\,905\,041$.

\begin{lemma}\label{lem:tail}
Assume (G) with $x_0 \le X_0$, and put $\varepsilon = \varepsilon(X_0)$. If
$\binomtwo{n} = B(a,k)$ with $p_a \ge X_0$, then
\[
\varepsilon\,(k-1)\lfloor k/2\rfloor \;\ge\; \ln(2 - 1/X_0).
\]
\end{lemma}
\begin{proof}
For $i \ge a$ we have $p_{i+1} \le p_i(1+\varepsilon(p_i)) \le p_i(1+\varepsilon)$, so
$r \le p(1+\varepsilon)^{k-1}$ and $(r/p)^{h} \le \exp(\varepsilon(k-1)h)$. Now apply
Lemma~\ref{lem:balance}, using $2 - 1/p \ge 2 - 1/X_0$.
\end{proof}

Evaluating Lemma~\ref{lem:tail} in high-precision arithmetic, a solution with
$p_a \ge X_0$ needs $k \ge 137$ under \cite{D10}, and $k \ge 1924$ under \cite{D18}.

\section{Computation}

\paragraph{Square test.} $B = n(n-1)/2$ if and only if $1 + 8B = (2n-1)^2$. When $B$ is
small we test this with an exact integer square root. For huge blocks we use $40$--$48$
filter primes $\ell$ and reject $B$ as soon as $1 + 8B \bmod \ell$ is a quadratic
non-residue (computed with the Jacobi-symbol algorithm). A true square is never rejected.
Every block that passes all the filters is re-checked exactly.

\paragraph{Search A (bounded $n$).} Fix $N$, $M = N(N-1)/2$ and $L = $ the bit length of $M$.
For a start $p$, let $e = \lfloor \log_2 p \rfloor$ and $K = \lceil L/e \rceil$. Blocks of
length $\ge K$ exceed $M$. For each start $p < X_0$ with $p \ge 41$, if
Lemma~\ref{lem:test} shows that Lemma~\ref{lem:balance} fails at length $K-1$, the start is
discarded (Corollary~\ref{cor:mono}). Otherwise all lengths $2 \le k < K$ are tested.
Starts $p < 41$ are always tested in full. Starts $p \ge X_0$ are handled by
Lemma~\ref{lem:tail} (the required $k$ exceeds $\ln M/\ln X_0$).

\paragraph{Search B (bounded $k$).} This is the same search with $K-1$ replaced by a fixed
maximal length $k_{\max}$.

\paragraph{Validation.}
\begin{itemize}
\item Three independent programs (a block search, a loop over $n$ with factorization, and
  Search A with and without the pruning of Lemma~\ref{lem:test}) agree for $N = 10^6$ and
  $10^9$.
\item In every run the number of starts processed equals $\pi(X_0) = 486\,570\,088$,
  computed independently.
\item No block has ever passed all the filters without being a genuine solution.
\end{itemize}

\begin{theorem}\label{thm:bounded}
If $4 \le n \le N$ and $\binomtwo{n}$ is a block, then $n \in \{4,6,15,21,715\}$, where
$N = 10^{12}$ unconditionally, $N = 10^{682}$ assuming \cite{D10}, and $N = 10^{9614}$
assuming \cite{D18}.
\end{theorem}

\begin{theorem}\label{thm:length}
If $\binomtwo{n}$ is a block of $k$ primes and $n \notin \{4,6,15,21,715\}$, then
$k \ge 137$ (assuming \cite{D10}) and $k \ge 1924$ (assuming \cite{D18}).
\end{theorem}
\begin{proof}
Starts $p_a \ge X_0$ are handled by Lemma~\ref{lem:tail}. Search~B with $k_{\max} = 136$
and $k_{\max} = 1923$ covers all starts $p_a < X_0$.
\end{proof}

\begin{theorem}\label{thm:small}
If $\binomtwo{n} = B(a,k)$ with $p_a < 100$ and $p_{a+k-1} \le 4\cdot10^9$, then
$n \in \{4,6,15,21,715\}$. In particular, $n(n-1) = r\#$ and $n(n-1) = 2\cdot r\#$ have no
other solutions with $r \le 4\cdot 10^9$.
\end{theorem}

\begin{table}[h]
\centering
\begin{tabular}{lrrr}
\toprule
run & time (8 cores) & starts enumerated & blocks tested \\
\midrule
Theorem~\ref{thm:bounded}, $N = 10^{12}$ (no pruning) & 160 s & --- & $2.7\cdot 10^{10}$ \\
Theorem~\ref{thm:bounded}, $N = 10^{682}$ & 33 s & 45\,555 & $5.1\cdot10^6$ \\
Theorem~\ref{thm:bounded}, $N = 10^{9614}$ & 480 s & 4\,702\,043 & $4.8\cdot10^9$ \\
Theorem~\ref{thm:length}, $k_{\max} = 1923$ & 138 s & 2\,896\,301 & $1.9\cdot10^9$ \\
Theorem~\ref{thm:small} & 775 s & 25 & $4.7\cdot10^9$ \\
\bottomrule
\end{tabular}
\caption{Computations (Apple M2).}
\end{table}

\section{Heuristics and obstacles}

The integers $n$ with $2B \mid n(n-1)$ are the $2^{k}$ residues obtained from the splittings
of the block's primes between $n$ and $n-1$. A solution is exactly the case where the
least nontrivial residue $n$ satisfies $\lambda := n(n-1)/(2B) = 1$. Modelling the residues
as $2^k$ random points in $[0, 2B)$ predicts $\lambda \approx B/4^k$. For blocks starting at
$2$ or $3$ and ending at $r \le 97$ (exhaustive enumeration), the observed minimum
$\lambda$ follows this prediction up to $\lambda \approx 10^{22}$. It equals $1$ only for
the known solutions.

Several natural routes to finiteness seem blocked:
\begin{itemize}
\item \textbf{abc.} The abc conjecture gives nothing, because $n(n-1)$ is squarefree up to
  the factor $2$. (Compare Brocard's problem, where abc does apply.)
\item \textbf{Size arguments.} The equation $|2\prod S - \prod T| = 1$ sits exactly at the
  trivial bound for a difference of distinct integers.
\item \textbf{Equidistribution.} The Fourier transform of the residue set is the Riesz
  product $\prod_q (1 + e(h c_q/q))$, with $c_q \equiv (2B/q)^{-1} \pmod q$. The set has
  only $2^k$ points, far too few for equidistribution at scale $\sqrt{B}$ to rule out a
  point in $[2, \sqrt{2B}]$.
\end{itemize}

\section{Formalization and data}

Lemma~\ref{lem:notprime} and the arithmetic core of Lemma~\ref{lem:balance} (the three
cases above, stated for lists of factors in $[p,r]$), together with the monotonicity in
$h$, are formalized in Lean~4 without Mathlib. The proofs use only the axioms
\texttt{propext} and \texttt{Quot.sound}. The code, logs and Lean sources are available at
[repository URL].

\paragraph{Acknowledgements.} The computations, proofs and formalization were developed with
substantial assistance from Claude (Anthropic), an AI model. The author checked [TO BE
COMPLETED BY AUTHOR: describe what you personally verified].

\begin{thebibliography}{9}
\bibitem{Ax} C.~Axler, New estimates for some functions defined over primes,
  arXiv:1703.08032.
\bibitem{D10} P.~Dusart, Estimates of some functions over primes without R.H.,
  arXiv:1002.0442 (2010).
\bibitem{D18} P.~Dusart, Explicit estimates of some functions over primes,
  \emph{Ramanujan J.} \textbf{45} (2018), 227--251, doi:10.1007/s11139-016-9839-4.
  [Cor.~5.5 checked only as quoted in \cite{Ax}; confirm against the published text.]
\bibitem{EG} P.~Erd\H{o}s and R.~L.~Graham, \emph{Old and New Problems and Results in
  Combinatorial Number Theory}, Monographies de L'Enseignement Math\'ematique 28, 1980.
\bibitem{EP} T.~Bloom, Erd\H{o}s Problems, Problem \#386,
  \url{https://www.erdosproblems.com/386}.
\bibitem{FC} Google DeepMind, formal-conjectures,
  \url{https://github.com/google-deepmind/formal-conjectures}.
\bibitem{NPP} C.~Nelson, D.~E.~Penney and C.~Pomerance, 714 and 715,
  \emph{J. Recreational Math.} \textbf{7}(2) (1974), 87--89.
\bibitem{OEIS} OEIS Foundation, Sequence A280992, \url{https://oeis.org/A280992}.
\end{thebibliography}

\end{document}
